% Глава 6. Ответы на RQ (после пилота), вклад, направления докторантуры (docs/roadmap.md, этап 2).
\chapter{Заключение}\label{ch:concl}

\section{Итоги}

Работа ставила целью спроектировать тренажёр чтения, следующий «Муаллим ас-Сани», и измерить, насколько
открытая модель распознавания фонем может выступать слушающим учителем для русскоязычных начинающих. Обзор
литературы показал, что звуковой пошаговый метод букваря подходит для автоматической обратной связи, что
явная обратная связь на уровне отдельных звуков --- условие, при котором обучение произношению работает
лучше всего, и что лучшая открытая модель для коранического арабского никогда не измерялась на букварном
чтении тех, для кого арабский не родной. Он показал и то, что ИИ-агенты склонны подгонять собственные тесты,
а письменные правила со временем слабеют, тогда как проверки, обеспечиваемые средой, помогают.

Созданный в ответ тренажёр превращает всю книгу в 125 уроков и 800 карточек, прослеживаемых до страницы и
строки; оценивает каждое голосовое на уровне фонем с полным следом решения; объясняет ошибки по-русски в
коротком формате, опирающемся на исследования обратной связи; и хранит собранные с согласия
псевдонимизированные данные, готовые к разметке учителями и оценке. Обвязка, в которой он создан,
обеспечивает замороженные тесты, защищённые зоны человека, контракты слоёв и зелёные проверки и ведёт
журналы, делающие разработку проверяемой.

\section{Ответы на исследовательские вопросы}

\paragraph{RQ-1} \pending{ответ после пилота}.
\paragraph{RQ-2} \pending{ответ после пилота}.
\paragraph{RQ-3} \pending{ответ после пилота}.
\paragraph{RQ-4} Собранные пока данные показывают, что обвязка заблокировала или сделала видимыми
классические сбои разработки с агентами --- подгонку тестов, размывание архитектуры, утечку секретов, --- а
несколько сбоев прошли все автоматические проверки и были пойманы только проверкой второго уровня или
человеком. \pending{количественный ответ в конце проекта}.

\section{Вклад}

\begin{enumerate}
  \item Протокол оценки открытой модели распознавания фонем Корана на букварном чтении начинающих
    относительно учителей, с предрегистрированными гипотезами и стандартными метриками MDD.
  \item Цифровая программа «Муаллим ас-Сани» в порядке книги, где каждая строка прослеживается до
    конкретного издания, и соответствие принципов Максуди функциям тренажёра.
  \item Тренажёр с открытой архитектурой, вердикты которого воспроизводимы и оспоримы.
  \item Задокументированная обвязка для создания исследовательского ПО ИИ-агентами с инцидентами,
    показывающими её ценность и слепые пятна.
\end{enumerate}

\section{Дальнейшая работа}

Магистерская работа задумана как первый этап докторского исследования. Модель данных уже хранит то, что
понадобится дальше: постоянные идентификаторы ученика, карточки и попытки; версии кода, программы, модели
и руководства по разметке в каждой записи; сырые выходы модели; жалобы; родной язык ученика. Отсюда шесть
направлений.
\begin{enumerate}
  \item \textbf{Модель для начинающих, изучающих арабский как второй язык.} Дообучить на речи начинающих,
    собранной с согласия, и измерить, уменьшатся ли ложные отклонения.
  \item \textbf{Калибровка и исход «не уверен».} Использовать уверенность по фонемам, чтобы выбирать пороги,
    при которых ложные тревоги редки, а пропуски безопасны, и отправлять неуверенные попытки учителю.
  \item \textbf{Адаптивное обучение.} Моделировать знание букв каждым учеником и планировать интервальное
    повторение; добавить упражнения на слух, которые, как показывает курс научного руководителя,
    укладываются в тот же порядок.
  \item \textbf{Правила таджвида.} Распространить проверку на \emph{ихфа}, \emph{идгам}, \emph{мадд} и
    \emph{калькаля}, где модель слабее всего.
  \item \textbf{Лонгитюдные и межъязыковые исследования.} Сравнить обучение с тренажёром, с учителем и
    смешанное; сравнить носителей русского, казахского, узбекского и татарского.
  \item \textbf{Открытый бенчмарк.} Опубликовать набор букварного чтения начинающих с согласием на
    публикацию --- бенчмарк, которого сейчас нет.
\end{enumerate}

\section{Заключительное замечание}

Максуди писал книгу для учеников, у которых был учитель. Эта работа не пытается убрать учителя; она
пытается дать ученикам внимательного слушателя на те часы, когда учителя рядом нет, и честно сказать,
насколько внимателен этот слушатель.
